\documentclass[11pt]{article}
\usepackage[margin=2.4cm]{geometry}
\usepackage{amsmath,amssymb,amsthm}
\usepackage[hidelinks]{hyperref}
\usepackage{enumitem}
% hidden links (hidelinks)

\newtheorem{theorem}{Theorem}[section]
\newtheorem{lemma}[theorem]{Lemma}
\newtheorem{proposition}[theorem]{Proposition}
\newtheorem{corollary}[theorem]{Corollary}
\newtheorem{definition}[theorem]{Definition}
\newtheorem{remark}[theorem]{Remark}

\title{Independent Reproof of the Mahar--Willner Doubling Lemmas\\ \large (Cell Extension and Zero-Truncation Induction, with the Boundary-Slope Transmission Coefficients)}
\author{BVE research project --- session 12}
\date{2026-08-05}

\begin{document}
\maketitle

\begin{abstract}
We give a fully self-contained proof of the two structural lemmas used by
Mahar--Willner \cite{mw} in their solution of the extremal problem for
$\lambda_2/\lambda_1$ of $-y''=\lambda\rho y$ (Dirichlet, $a\le\rho\le 1$):
Lemma 1 (the \emph{cell-extension} identity
$\lambda_{2n}(\rho_n)/\lambda_n(\rho_n)=\lambda_2(\rho_0)/\lambda_1(\rho_0)$
for the $1/n$-periodic ``blow-up'' $\rho_n$ of an arbitrary $\rho_0$) and
Lemma 2 (the \emph{zero-truncation} induction establishing
$\mu_{2n,n}(a)=\mu(a)$, $\nu_{2n,n}(a)=\nu(a)$).  The only external facts used
are the classical Sturm oscillation and comparison theorems.  In the minimum
case we give a variant of the truncation argument which is cleaner than the
``similarly'' remark in \cite{mw}: the side on which one truncates and the
direction of the shooting are dictated by the position of $z_0$ relative to
$z_{2k}$, and the induction hypothesis is needed only in one of the two cases.
The proof makes the sup-ratio theorem
$\sup_{n,\rho}\lambda_{n+1}/\lambda_n=\nu(R)$ of the accompanying report
\emph{fully self-contained}.
\end{abstract}

\section{Setting and notation}

Throughout, $\rho$ denotes a piecewise constant density on
$[-1/2,1/2]$ with $a\le\rho\le 1$ pointwise ($0<a<1$), and
$\mathcal P_a$ is the class of such densities.  Consider
\begin{equation}
	-y''=\lambda\rho(x)\,y,\qquad y(-1/2)=y(1/2)=0,
	\label{eq:main}
\end{equation}
with eigenvalues $0<\lambda_1(\rho)<\lambda_2(\rho)<\cdots$ (simplicity and
discreteness are classical).  Write
\[
	\mu(a):=\inf_{\rho\in\mathcal P_a}\frac{\lambda_2}{\lambda_1},\qquad
	\nu(a):=\sup_{\rho\in\mathcal P_a}\frac{\lambda_2}{\lambda_1},
	\qquad
	\mu_{2n,n}(a):=\inf_{\rho\in\mathcal P_a}\frac{\lambda_{2n}}{\lambda_n},
	\qquad
	\nu_{2n,n}(a):=\sup_{\rho\in\mathcal P_a}\frac{\lambda_{2n}}{\lambda_n}.
\]
The eigenfunction of $\lambda_j$ is denoted $y_j$; it is normalized
arbitrarily, and we always take $y_j'(-1/2)\neq 0$ (true up to rescaling).

\medskip
\noindent\textbf{Preliminary facts.}  We shall use the following classical
results; proofs are sketched only where the formulation is non-standard.

\begin{proposition}[affine invariance of ratios]\label{prop:affine}
	Let $L(x)=cx+d$ with $c\neq 0$ map $[x_0,x_1]$ onto $[-1/2,1/2]$.
	If $\lambda_1<\lambda_2$ are the first two Dirichlet eigenvalues of
	\eqref{eq:main} on $[L(x_0),L(x_1)]$ with density $\rho$, then the first
	two Dirichlet eigenvalues of $y''+\lambda\rho(L(x))y=0$ on $[x_0,x_1]$
	are $\lambda_1/c^2$, $\lambda_2/c^2$.  In particular the ratio
	$\lambda_2/\lambda_1$ is invariant.
\end{proposition}

\begin{proposition}[Sturm oscillation]\label{prop:osc}
	The $j$-th eigenfunction $y_j$ of \eqref{eq:main} has exactly $j-1$
	(interior) zeros in $(-1/2,1/2)$, all simple.
\end{proposition}

\begin{proposition}[Sturm interlacing]\label{prop:interl}
	If $j<k$ and $z_1<\cdots<z_{j-1}$ are the zeros of $y_j$, then $y_k$ has
	at least one zero in each of the $j$ intervals
	$(-1/2,z_1),(z_1,z_2),\ldots,(z_{j-1},1/2)$.
\end{proposition}

The backbone of the truncation argument is the following elementary
monotonicity of the zeros of an initial-value problem (IVP).

\begin{proposition}[zero motion]\label{prop:zeromotion}
	Fix $c\neq 0$ and let $y(\cdot;\lambda)$ be the solution of
	\begin{equation*}
		y''+\lambda\rho(x)y=0,\qquad y(-1/2;\lambda)=0,\quad
		y'(-1/2;\lambda)=c.
	\end{equation*}
	For $j\ge 1$ let $z_j(\lambda)$ be the $j$-th zero of $y(\cdot;\lambda)$
	to the right of $-1/2$, whenever it exists.  Then each $z_j$ is
	continuous and \emph{strictly decreasing} in $\lambda$.  The analogous
	statement holds for the IVP posed at $x=1/2$: the $j$-th zero to the
	left of $1/2$ is strictly increasing in $\lambda$.
\end{proposition}

\begin{proof}
	Continuity is standard (smooth dependence on the parameter $\lambda$
	for ODEs with simple zeros).  For the monotonicity, use the Prufer
	representation $y=r\sin\theta$, $y'=\sqrt{\lambda\rho}\,r\cos\theta$,
	$\theta(-1/2)=0$.  Then $\theta$ satisfies
	$\theta'=\sqrt{\lambda\rho}\cos^2\theta+\frac{(\sqrt{\lambda\rho})'}{}
	\cdot$ (the exact form is immaterial); the comparison theorem for the
	Prufer angle gives $\theta(x;\lambda_1)<\theta(x;\lambda_2)$ for
	$\lambda_1<\lambda_2$ and every fixed $x$ (strictly, since
	$\sqrt{\lambda_2\rho}-\sqrt{\lambda_1\rho}>0$ a.e.).  The zeros of
	$y(\cdot;\lambda)$ are the solutions of $\theta(x;\lambda)=j\pi$, so
	$z_j(\lambda_2)<z_j(\lambda_1)$.  The statement at $x=1/2$ follows by
	the reflection $x\mapsto 1/2-x$.
\end{proof}

\noindent An equivalent formulation, used repeatedly below, is:

\begin{corollary}\label{cor:crossing}
	Let $\bar y(\cdot;\lambda)$ be an IVP solution as above and let
	$x_*\in(-1/2,1/2)$ be a fixed point which is not a zero of
	$\bar y(\cdot;\lambda_0)$.  As $\lambda$ increases (resp.\ decreases)
	from $\lambda_0$, the value $\bar y(x_*;\lambda)$ changes sign exactly
	when a zero of $\bar y$ crosses $x_*$; zeros cross one at a time, and
	the count of zeros of $\bar y$ strictly inside a fixed interval whose
	boundary is $x_*$ changes by $\pm 1$ at each crossing.
\end{corollary}

\section{Lemma 1: the cell extension}\label{sec:lemma1}

\begin{lemma}[cell extension]\label{lem:cell}
	For every $\rho_0\in\mathcal P_a$, every integer $n\ge 1$, and the
	$1/n$-periodic density
	\begin{equation*}
		\rho_n(x):=\rho_0\bigl(n(x+\tfrac12)-\tfrac12\bigr),
		\qquad -\tfrac12\le x\le -\tfrac12+\tfrac1n,
		\quad\text{extended with period }1/n,
	\end{equation*}
	one has $\rho_n\in\mathcal P_a$ and
	\begin{equation}
		\frac{\lambda_{2n}(\rho_n)}{\lambda_n(\rho_n)}
		=\frac{\lambda_2(\rho_0)}{\lambda_1(\rho_0)}.
		\label{eq:cellidentity}
	\end{equation}
\end{lemma}

\begin{proof}
	Fix $\rho_0$ and write $y_1,y_2$ for its first two eigenfunctions, with
	eigenvalues $\lambda_1,\lambda_2$.  Since $y_1>0$ on $(-1/2,1/2)$,
	$y_1(\pm 1/2)=0$, we have $y_1'(-1/2)>0>y_1'(1/2)$; define the
	transmission coefficient
	\begin{equation*}
		s:=-\frac{y_1'(1/2)}{y_1'(-1/2)}>0.
	\end{equation*}
	Similarly $y_2$ has exactly one interior zero, so its boundary slopes
	have opposite signs and
	$t:=-y_2'(1/2)/y_2'(-1/2)>0$ satisfies
	$y_2'(1/2)=-t\,y_2'(-1/2)$.

	Set $x_j:=-1/2+j/n$ and, for $x\in[x_j,x_{j+1}]$,
	$u_j(x):=n(x-x_j)-1/2\in[-1/2,1/2]$.  Define
	\begin{equation*}
		\widetilde y(x):=(-s)^j\,y_1(u_j(x)),\qquad
		\widetilde f(x):=(-t)^j\,y_2(u_j(x)),
		\qquad x\in[x_j,x_{j+1}],\ j=0,\ldots,n-1.
	\end{equation*}
	Both are $C^1$ across every cell boundary $x_{j+1}$: continuity holds
	because $y_1(\pm 1/2)=y_2(\pm 1/2)=0$, and the slopes match because
	\begin{equation*}
		(-s)^j y_1'(1/2)=(-s)^{j+1}y_1'(-1/2),\qquad
		(-t)^j y_2'(1/2)=(-t)^{j+1}y_2'(-1/2).
	\end{equation*}
	On cell $j$ one has $\widetilde y''=-n^2\lambda_1\rho_n\widetilde y$ and
	$\widetilde f''=-n^2\lambda_2\rho_n\widetilde f$, so both are
	eigenfunctions of \eqref{eq:main} with $\rho=\rho_n$, of eigenvalues
	$n^2\lambda_1$ and $n^2\lambda_2$.  Their boundary values at
	$\pm 1/2$ vanish.

	Zeros: $\widetilde y$ vanishes exactly at the cell boundaries, i.e.\ at
	$n-1$ interior points, hence it is the $n$-th eigenfunction
	(Proposition~\ref{prop:osc}) and
	$\lambda_n(\rho_n)=n^2\lambda_1(\rho_0)$.  The function
	$\widetilde f$ vanishes at the $n-1$ cell boundaries and at one point
	inside each cell (the image of the interior zero of $y_2$), i.e.\ at
	$2n-1$ interior points, hence it is the $2n$-th eigenfunction and
	$\lambda_{2n}(\rho_n)=n^2\lambda_2(\rho_0)$.  Dividing gives
	\eqref{eq:cellidentity}.
\end{proof}

\begin{corollary}\label{cor:lemma1}
	$\mu_{2n,n}(a)\le\mu(a)$ and $\nu_{2n,n}(a)\ge\nu(a)$ for all $n\ge 1$.
\end{corollary}

\begin{proof}
	For any $\rho_0$, \eqref{eq:cellidentity} says that the value
	$\lambda_2(\rho_0)/\lambda_1(\rho_0)$ is attained by the ratio
	$\lambda_{2n}/\lambda_n$ at $\rho_n\in\mathcal P_a$.  Taking the
	supremum (resp.\ infimum) over $\rho_0$ gives
	$\nu_{2n,n}(a)\ge\nu(a)$ (resp.\ $\mu_{2n,n}(a)\le\mu(a)$).
\end{proof}

\section{Lemma 2: the zero-truncation induction}\label{sec:lemma2}

Throughout this section, $k\ge 1$, $m:=k+1$, and $\rho$ is an arbitrary
element of $\mathcal P_a$.  We write
\begin{equation*}
	z_0:=\text{rightmost interior zero of }y_{m},\qquad
	z_1:=\text{second-from-right interior zero of }y_{2m},
\end{equation*}
and $z_1<z_2<\cdots<z_{2k+1}$ for all interior zeros of $y_{2m}$ (so
$z_{2k}=z_1$).  By Proposition~\ref{prop:interl}, $y_{2m}$ has at least one
zero in each of the $k$ intervals between consecutive zeros of $y_m$ inside
$(-1/2,z_0]$; in particular at least $k$ zeros of $y_{2m}$ lie strictly
inside $(-1/2,z_0)$ and at least one lies in $(z_0,1/2)$.

\subsection{The maximum case}

\begin{lemma}[truncation, max]\label{lem:max}
	For every $\rho\in\mathcal P_a$ and every $k\ge 1$,
	\begin{equation*}
		\frac{\lambda_{2(k+1)}(\rho)}{\lambda_{k+1}(\rho)}
		\le\max\bigl\{\nu_{2k,k}(a),\,\nu(a)\bigr\}.
	\end{equation*}
\end{lemma}

\begin{proof}
	We fix $\rho$ and suppress it from the notation.
	\emph{Case A: $z_0\le z_1$ (i.e.\ $z_0\le z_{2k}$).}
	Let $\bar y(\cdot;\lambda)$ be the IVP solution
	$y''+\lambda\rho y=0$, $y(-1/2;\lambda)=0$, $y'(-1/2;\lambda)=1$, and let
	$j$ be the number of zeros of $y_{2m}$ strictly inside $(-1/2,z_0)$.
	Since $z_{2k}\ge z_0$, we have $j\le 2k-1$.  By
	Proposition~\ref{prop:zeromotion} and Corollary~\ref{cor:crossing},
	there is a smallest $\bar\lambda\ge\lambda_{2m}$ with
	$\bar y(z_0;\bar\lambda)=0$; at this value the zero of $\bar y$ lying
	just to the right of $z_0$ at $\lambda=\lambda_{2m}$ has moved to $z_0$,
	so $\bar y(\cdot;\bar\lambda)$ has exactly $j\le 2k-1$ zeros strictly
	inside $(-1/2,z_0)$.

	Consider the truncated problem
	$y''+\lambda\rho y=0$ on $(-1/2,z_0)$ with
	$y(-1/2)=y(z_0)=0$.  Its $k$-th eigenfunction is $y_m$ restricted
	($k-1$ zeros inside), with eigenvalue $\lambda_{k+1}$; its
	$(j+1)$-st eigenfunction is $\bar y(\cdot;\bar\lambda)$ ($j$ zeros
	inside), with eigenvalue $\bar\lambda$.  Since $j\ge k$
	(Proposition~\ref{prop:interl}), $k<j+1\le 2k$, and
	$\lambda_{j+1}\le\lambda_{2k}$ pointwise, the affine
	normalization (Proposition~\ref{prop:affine}) gives
	\begin{equation*}
		\frac{\bar\lambda}{\lambda_{k+1}}
		\le\sup_{\rho'\in\mathcal P_a}\frac{\lambda_{j+1}(\rho')}
		{\lambda_k(\rho')}
		\le\sup_{\rho'\in\mathcal P_a}\frac{\lambda_{2k}(\rho')}
		{\lambda_k(\rho')}
		=\nu_{2k,k}(a).
	\end{equation*}
	Since $\lambda_{2m}\le\bar\lambda$, we obtain
	$\lambda_{2m}/\lambda_{k+1}\le\nu_{2k,k}(a)$.

	\emph{Case B: $z_1<z_0$ (i.e.\ $z_{2k}<z_0$).}
	Now $y_{2m}$ has exactly one zero, $z_{2k+1}$, strictly inside
	$(z_0,1/2)$.  Let $\bar y(\cdot;\lambda)$ be the IVP solution posed at
	$x=1/2$: $y(1/2;\lambda)=0$, $y'(1/2;\lambda)=1$.  As $\lambda$
	increases, its zeros move toward $1/2$; there is a smallest
	$\bar\lambda\ge\lambda_{2m}$ with $\bar y(z_0;\bar\lambda)=0$, and at
	this value $\bar y$ has exactly one zero strictly inside $(z_0,1/2)$
	(namely the image of $z_{2k+1}$).  Hence $\bar\lambda$ is the second
	eigenvalue of the truncated problem on $(z_0,1/2)$ whose first
	eigenfunction is $y_m$ restricted (no zeros inside), of eigenvalue
	$\lambda_{k+1}$.  Normalizing affinely,
	\begin{equation*}
		\frac{\bar\lambda}{\lambda_{k+1}}\le\nu(a),
		\qquad\text{so}\qquad
		\frac{\lambda_{2m}}{\lambda_{k+1}}\le\nu(a).
	\end{equation*}
	In both cases $\lambda_{2m}/\lambda_{k+1}\le
	\max\{\nu_{2k,k}(a),\nu(a)\}$, as claimed.
\end{proof}

\begin{theorem}[MW Lemma 2, max]\label{thm:max}
	$\nu_{2n,n}(a)=\nu(a)$ for all $n\ge 1$.
\end{theorem}

\begin{proof}
	Induction on $k$ with $n=k+1$.  The case $k=1$ is
	$\nu_{2,1}(a)=\nu(a)$ by definition.  If
	$\nu_{2k,k}(a)\le\nu(a)$ (indeed equality holds by
	Corollary~\ref{cor:lemma1} and the induction hypothesis), then
	Lemma~\ref{lem:max} applied to an arbitrary $\rho$ gives
	$\lambda_{2(k+1)}/\lambda_{k+1}\le\nu(a)$; taking the supremum over
	$\rho$ yields $\nu_{2(k+1),(k+1)}(a)\le\nu(a)$.  The reverse inequality
	is Corollary~\ref{cor:lemma1}.
\end{proof}

\subsection{The minimum case}

For the minimum we use the following variant, which is the precise form of
the ``minimization following similarly'' in \cite{mw}.  Let
$c_0$ be the number of zeros of $y_{2m}$ strictly inside $(z_0,1/2)$.
By interlacing $c_0\ge 1$.

\begin{lemma}[truncation, min]\label{lem:min}
	For every $\rho\in\mathcal P_a$ and every $k\ge 1$,
	\begin{equation*}
		\frac{\lambda_{2(k+1)}(\rho)}{\lambda_{k+1}(\rho)}
		\ge\min\bigl\{\mu_{2k,k}(a),\,\mu(a)\bigr\}.
	\end{equation*}
\end{lemma}

\begin{proof}
	Fix $\rho$.
	\emph{Case A: $c_0\ge 2$ (i.e.\ $z_0<z_{2k}$).}
	Consider the truncated problem on $(z_0,1/2)$; its first eigenfunction
	is $y_m$ restricted, of eigenvalue $\lambda_{k+1}$.  Let
	$\bar y(\cdot;\lambda)$ be the IVP solution posed at $1/2$ with initial
	slope $1$.  As $\lambda$ decreases from $\lambda_{2m}$, the zeros of
	$\bar y$ move to the left (Proposition~\ref{prop:zeromotion}) and exit
	the interval $(z_0,1/2)$ through $z_0$ one at a time
	(Corollary~\ref{cor:crossing}); the count strictly inside
	$(z_0,1/2)$ starts at $c_0\ge 2$ and decreases by $1$ at each crossing.
	Let $\bar\lambda$ be the value at which the count drops from $2$ to $1$,
	i.e.\ the largest $\lambda\le\lambda_{2m}$ with $\bar y(z_0;\lambda)=0$
	after exactly $c_0-2$ crossings.  Then $\bar\lambda\le\lambda_{2m}$,
	$\bar y(\cdot;\bar\lambda)$ has exactly one zero strictly inside
	$(z_0,1/2)$, and it is an eigenfunction of the truncated problem; hence
	$\bar\lambda$ is its \emph{second} eigenvalue.  Normalizing affinely and
	using the definition of $\mu$,
	\begin{equation*}
		\frac{\bar\lambda}{\lambda_{k+1}}\ge\mu(a),
		\qquad\text{so}\qquad
		\frac{\lambda_{2m}}{\lambda_{k+1}}\ge\mu(a).
	\end{equation*}

	\emph{Case B: $c_0=1$ (i.e.\ $z_{2k}<z_0$).}
	Then all $2k$ zeros $z_1,\ldots,z_{2k}$ of $y_{2m}$ lie strictly inside
	$(-1/2,z_0)$.  Let $\bar y(\cdot;\lambda)$ be the IVP solution posed at
	$-1/2$ with initial slope $1$.  As $\lambda$ decreases, zeros move to
	the right; let $\bar\lambda$ be the largest $\lambda\le\lambda_{2m}$
	such that $\bar y(z_0;\lambda)=0$ (the first crossing).  Then
	$\bar y(\cdot;\bar\lambda)$ has exactly $2k-1$ zeros strictly inside
	$(-1/2,z_0)$, so it is the $2k$-th eigenfunction of the truncated
	problem on $(-1/2,z_0)$, while $y_m$ restricted is its $k$-th
	eigenfunction with eigenvalue $\lambda_{k+1}$.  Hence, after affine
	normalization,
	\begin{equation*}
		\frac{\bar\lambda}{\lambda_{k+1}}\ge\mu_{2k,k}(a),
		\qquad\text{so}\qquad
		\frac{\lambda_{2m}}{\lambda_{k+1}}\ge\mu_{2k,k}(a).
	\end{equation*}
	In both cases the desired lower bound holds.
\end{proof}

\begin{theorem}[MW Lemma 2, min]\label{thm:min}
	$\mu_{2n,n}(a)=\mu(a)$ for all $n\ge 1$.
\end{theorem}

\begin{proof}
	Induction on $k$ with $n=k+1$.  Base case $k=1$:
	$\mu_{2,1}(a)=\mu(a)$ by definition.  If
	$\mu_{2k,k}(a)\ge\mu(a)$, then Lemma~\ref{lem:min} applied to an
	arbitrary $\rho$ gives $\lambda_{2(k+1)}/\lambda_{k+1}\ge\mu(a)$;
	taking the infimum yields
	$\mu_{2(k+1),(k+1)}(a)\ge\mu(a)$.  The reverse inequality is
	Corollary~\ref{cor:lemma1}.
\end{proof}

\section{Consequences}

\begin{theorem}[MW Theorem 3]\label{thm:theorem3}
	If $\rho_0$ extremizes $\lambda_2/\lambda_1$ in $\mathcal P_a$, then the
	cell extension $\rho_n$ of Lemma~\ref{lem:cell} extremizes
	$\lambda_{2n}/\lambda_n$.  Consequently
	\begin{equation*}
		\mu_{2n,n}(a)=\mu(a),\qquad
		\nu_{2n,n}(a)=\nu(a)\qquad(n\ge 1).
	\end{equation*}
\end{theorem}

\begin{proof}
	If $\rho_0$ maximizes, then $\nu(a)=\lambda_2(\rho_0)/\lambda_1(\rho_0)
	=\lambda_{2n}(\rho_n)/\lambda_n(\rho_n)\le\nu_{2n,n}(a)$ by definition,
	and Theorem~\ref{thm:max} gives the reverse inequality; hence $\rho_n$
	maximizes $\lambda_{2n}/\lambda_n$.  The minimum case is analogous with
	Theorem~\ref{thm:min}.
\end{proof}

\begin{corollary}[self-contained sup-ratio theorem]\label{cor:sup}
	Let $R>1$ and $0<a\le\rho\le A=aR$.  For the problem
	$-y''=\lambda\rho y$ (Dirichlet on $[-1/2,1/2]$),
	\begin{equation*}
		\sup_{n\ge 1}\sup_{\rho}\frac{\lambda_{n+1}}{\lambda_n}
		=\nu(R)=\Bigl(\frac{\arccos(-\sqrt{R}/(\sqrt{R}+1))}
		{\arccos(\sqrt{R}/(\sqrt{R}+1))}\Bigr)^2,
	\end{equation*}
	attained at $n=1$ and the symmetric two-jump configuration
	$[a,A,a]$.
\end{corollary}

\begin{proof}
	By scaling $\rho\mapsto\rho/a$ we may assume $a=1$, $A=R$.
	The upper bound: $\lambda_{n+1}\le\lambda_{2n}$ and
	Theorem~\ref{thm:max} give
	$\lambda_{n+1}/\lambda_n\le\nu_{2n,n}(R)=\nu(R)$.
	The lower bound and the closed formula for $\nu(R)$, together with the
	explicit extremizer $[1,R,1]$ with block widths
	$s t,t,s t$, $s=\sqrt R$, $t=1/(2s+1)$, are established in the
	accompanying report \texttt{SL\_ratio\_proof.pdf} (balanced-phase
	computation).  This proves the claim without appealing to
	\cite[Lemmas 1--2]{mw}.
\end{proof}

\section{Numerical audit}

All computations use exact piecewise-constant transfer-matrix propagation
(double precision) with bisection for eigenvalues and zero locations.

\begin{enumerate}[leftmargin=1.6em]
	\item \emph{Lemma~\ref{lem:cell}.}  For 8 random densities
	$\rho_0$ with 2--4 jumps and $1\le\rho_0\le 4$, the identity
	\eqref{eq:cellidentity} was verified for $n=2,3$ to relative error
	$\le 4.5\times10^{-16}$.  Additionally the identities
	$\lambda_n(\rho_n)=n^2\lambda_1(\rho_0)$,
	$\lambda_{2n}(\rho_n)=n^2\lambda_2(\rho_0)$ were checked to 6 digits,
	confirming the $C^1$ pasting structure.  (Script
	\texttt{op05\_mw\_lemma12\_audit.py}.)
	\item \emph{Lemma~\ref{lem:max}.}  For 24 random configurations
	(10 case A, 14 case B, $k=1,2,3$), the shooting value
	$\bar\lambda$ was computed by sign-change bisection and verified to
	satisfy: case A --- $\bar\lambda$ is the $(j+1)$-st eigenvalue of the
	truncated problem with $j\le 2k-1$ zeros, $j\ge k$, and
	$\bar\lambda\ge\lambda_{2m}$; case B --- $\bar\lambda$ is the second
	eigenvalue of the $(z_0,1/2)$-problem, with exactly one interior zero,
	and $\bar\lambda\ge\lambda_{2m}$.  All checks passed to relative error
	$<10^{-4}$ (limited by the zero-location precision of $z_0$).
	(Script \texttt{op05\_mw\_lemma2\_audit.py}.)
	\item \emph{Lemma~\ref{lem:min}.}  For 20 random configurations
	(14 case A, 6 case B, $k=1,2,3$), case A: the shooting value equals
	the second eigenvalue of $(z_0,1/2)$ and $\bar\lambda\le\lambda_{2m}$
	(with ratios $e_2/\lambda_m\in[3.65,4.18]$);
	case B: $\bar\lambda$ equals the $(2k)$-th eigenvalue of
	$(-1/2,z_0)$ and $\bar\lambda\le\lambda_{2m}$.  All checks passed.
	(Scripts \texttt{op05\_mw\_mincase\_audit2.py}.)
	\item \emph{Zero motion (Proposition~\ref{prop:zeromotion}).}
	For 5 random densities, the ordered zero positions of the IVP solution
	at $\lambda$ and $1.05\lambda$ are strictly ordered as predicted.
	(Script \texttt{op05\_mw\_extra\_checks.py}.)
\end{enumerate}

\section{Audit notes}

\begin{itemize}
	\item The case split in Lemma~\ref{lem:max} is exhaustive:
	$z_0\le z_1$ or $z_1<z_0$.  The degenerate case $z_0=z_1$ is covered by
	case A with $j\le 2k-1$ (then $z_{2k}=z_0$ and $\bar\lambda=\lambda_{2m}$
	is admissible).
	\item The identity $j\ge k$ in case A of Lemma~\ref{lem:max} follows
	from interlacing: each of the $k$ intervals between $-1/2$ and the
	successive zeros of $y_m$ contains a zero of $y_{2m}$ strictly inside
	$(-1/2,z_0)$.
	\item In case A of Lemma~\ref{lem:min} the ``count drops from 2 to 1''
	uses the monotonicity of the zero count in $\lambda$; the value
	$\bar\lambda$ is obtained after exactly $c_0-2$ crossings, all of which
	occur at $\lambda<\lambda_{2m}$, so $\bar\lambda\le\lambda_{2m}$.
	\item We do not need the existence of extremizers for the inequalities
	$\nu_{2n,n}\le\nu$ and $\mu_{2n,n}\ge\mu$: Lemma~\ref{lem:max}
	(resp.\ Lemma~\ref{lem:min}) is proved for an \emph{arbitrary}
	$\rho\in\mathcal P_a$, and only the final supremum (resp.\ infimum) is
	taken.  This removes the compactness question entirely.
	\item The transmission coefficients $s,t$ are only used in
	Lemma~\ref{lem:cell}; the equality $s=t$ is \emph{not} required for the
	doubling identities (it is used in \cite{mw} only to prove the
	symmetry of the extremizer, Theorem 2 there).
	\item All external facts (oscillation, interlacing, Prufer comparison)
	are classical; the remaining steps are elementary.
\end{itemize}

\begin{thebibliography}{9}
\bibitem{mw} T.~J.~Mahar, B.~E.~Willner, \emph{An extremal eigenvalue
	problem}, Comm. Pure Appl. Math. 29 (1976), no.~5, 517--529.
	\url{https://doi.org/10.1002/cpa.3160290505}
\bibitem{keller} J.~B.~Keller, \emph{The minimum ratio of two eigenvalues},
	SIAM J. Appl. Math. 31 (1976), no.~3, 485--491.
	\url{https://doi.org/10.1137/0131042}
\end{thebibliography}

\end{document}

